\section{Розділ~\ref{sec:recognition}. Розпізнавання}

Швидкість розпізнавання, будову перших ступенів, лінійне читання відповіді й навчання на неперервності часу виміряно на людині й мавпі; зведення всього ланцюжка до двох операцій і навчання на відео перевірено на моделях книги, а пояснення ілюзій --- від добре обґрунтованих до спірних.

{\footnotesize
\setlength{\tabcolsep}{3pt}\renewcommand{\arraystretch}{1.15}
\begin{longtable}{>{\raggedright}p{0.5cm}>{\raggedright}p{4.0cm}>{\raggedright}p{5.6cm}>{\raggedright}p{2.3cm}>{\raggedright\arraybackslash}p{1.7cm}}
\toprule
№ & Твердження & Дослід і що виміряно & Джерело & Статус \\
\midrule\endfirsthead
\toprule
№ & Твердження & Дослід і що виміряно & Джерело & Статус \\
\midrule\endhead
1 & Порівняння зі зразком за пікселями при зсуві вгадує не краще за монетку; пара ступенів~\eqref{eq:scstage} упізнає фігуру в будь-якому місці & \texttt{models/\allowbreak recognition\_models.py}, «сітківка» з $24$ точок, $4000$ проб: зразок --- $0.49$, $0.51$, $0.52$ за частки перевернутих точок $0$, $0.1$, $0.2$; пара $S$ і $C$ --- $1.00$, $0.86$, $0.70$ & виведення в розділі & модель \\
2 & Тварину на картинці впізнають за $\sim150$\,мс, один прохід десятком ступенів по $10$--$20$\,мс & Людина, задача «є тварина чи ні» на фотографіях, показаних на $20$\,мс: викликані потенціали для двох відповідей розходяться через $\sim150$\,мс. Мавпа: час першої відповіді зростає від ступеня до ступеня (V1 раніше за всіх, далі V2 і V4). Кількість ступенів і «нуль або один імпульс на ступінь» --- висновок із цих чисел & \cite{Thorpe1996,Schmolesky1998} & виміряно \\
3 & Ступінь $S$ --- І за частинами, ступінь $C$ --- найбільше за положеннями (твердження~\ref{prop:conveyor}) & Кішка, первинна зорова кора: прості клітини відповідають на край певної орієнтації в певному місці, складні --- на ту саму орієнтацію в більшій ділянці. Внутрішньоклітинний запис складних клітин: відповідь на дві смужки в багатьох клітин близька до більшої з двох відповідей, у середньому найближча до операції max. Найбільше через взаємне гальмування --- твердження~\ref{prop:wta}, ділення на загальну активність --- огляд & \cite{HubelWiesel1962,Lampl2004,Carandini2012} & виміряно \\
4 & Вище ланцюжком поєднання більші й складніші, відповідь дедалі менше залежить від положення & Мавпа, спрощення стимулів до «критичної ознаки»: клітини V4 і задньої нижньоскроневої кори відповідають на складні поєднання форми, кольору й текстури за малого поля; у передній нижньоскроневій корі поле більше. Чергування $S$ і $C$ у моделі відтворює цю картину & \cite{KobatakeTanaka1994,Riesenhuber1999} & виміряно \\
5 & Згорткові мережі повторюють це чергування й найкраще передбачають відповіді верхніх ступенів; предмет читається лінійно за частотами групи нейронів & Мережа, навчена розпізнавати, без підгонки до мозку: верхній шар передбачає відповіді нейронів нижньоскроневої кори мавпи, середні --- відповіді V4. Лінійний класифікатор за активністю $\sim100$ випадково вибраних клітин у вікнах від $12.5$\,мс упізнає предмет і категорію за різних положення й розміру & \cite{Fukushima1980,Yamins2014,Hung2005} & виміряно \\
6 & Правило Гебба зі слідом~\eqref{eq:traceinvariance} навчає інваріантності на відео без підписів, якщо слід не коротший за $\sim20$\,мс & Ідея повільних ознак і правила зі слідом --- теорія. \texttt{models/\allowbreak recognition\_models.py}, два $C$-нейрони, $16$ входів, $10$ прогонів: пауза між предметами $0.3$\,с. При $\tau_{\text{сл}}=5$\,мс збіг із фігурою $0.52$ у середньому, при $10$\,мс --- $0.56$; при $20$, $50$ і $200$\,мс --- $1.00$ в усіх прогонах (при $30$\,мс один прогін із десяти помиляється) & \cite{Foldiak1991,WiskottSejnowski2002} & модель \\
7 & У мозку інваріантність навчається на неперервності часу & Мавпа: предмет непомітно підміняли під час стрибка ока в певне місце поля; позиційна інваріантність нейронів нижньоскроневої кори змінювалася відповідно до підміни, значуще вже за годину. Що це головне правило навчання зору --- гіпотеза & \cite{LiDiCarlo2008} & виміряно \\
8 & Рух: детектори напрямку, далі та сама пара І/АБО за великими ділянками & Детектори напрямку в сітківці кроля; в ділянці MT мавпи всі $110$ записаних нейронів вибіркові до напрямку, відповідь на рух сильніша, ніж у V1 & \cite{Barlow1965,Albright1984} & виміряно \\
9 & Верхні ступені надсилають униз передбачення, угору йде помилка & Модель пояснює позакласичні ефекти полів первинної зорової кори; окремі спостереження узгоджуються (огляд), як загальну будову конвеєра не доведено & \cite{RaoBallard1999,Keller2018} & гіпотеза \\
10 & Складні картинки потребують зворотних зв'язків і кількох кіл & Мавпа: для «складних» картинок рішення в нижньоскроневій корі складається на $\sim30$\,мс пізніше; ці пізні відповіді погано передбачає пряма мережа й краще --- мережі зі зворотними зв'язками або глибші & \cite{Kar2019} & виміряно \\
11 & Непомітна підробка пікселів обманює мережу; зір людини набагато стійкіший, але не невразливий & На мережах: мале спеціально дібране збурення змінює відповідь; приклад «панда $\to$ гібон» --- із другої праці. У людини за короткого показу підробки, що добре переносяться між мережами, зсувають вибір & \cite{Szegedy2014,Goodfellow2015,Elsayed2018} & виміряно \\
12 & Ілюзії очікування: ненадійний вхід поступається очікуванню --- бліде рухається «повільніше», сукню бачать по-різному (підрозділ~\ref{sec:illusions}) & Ґратки меншого контрасту здаються рухомими повільніше. Опитування $1401$ людини: сукню бачать у трьох стійких варіантах, підказка про освітлення перемикає колір; у $\sim13\,000$ спостерігачів вирішує припущення про освітлення. Баєсова модель з очікуванням повільних рухів пояснює низку ілюзій руху & \cite{Thompson1982,LaferSousa2015,Wallisch2017,Weiss2002,Purves2011} & гіпотеза \\
13 & Регулювання підсилення: водоспад, нахил і контраст; ґратка Германа --- не гальмування сусідів & Детектори напрямку сітківки кроля за тривалого руху знижують частоту, а після зупинки падають нижче за фон. Ілюзію нахилу відтворює модель із діленням на активність сусідів. Зміни ґратки, що не змінюють відповіді вихідних нейронів сітківки, прибирають плями --- пояснення сусідським гальмуванням відкинуто & \cite{BarlowHill1963,Schwartz2009,SchillerCarvey2005} & виміряно \\
14 & Компенсація затримок: рухомий предмет здається попереду спалаху & Ілюзію виміряно. Психофізика суперечить і продовженню руху вперед, і різниці затримок; запропоновано пояснення заднім числом --- за подіями $\sim80$\,мс після спалаху. Суперечку не розв'язано & \cite{Nijhawan1994,EaglemanSejnowski2000} & гіпотеза \\
15 & Нерухомі «змії» обертаються: різниця затримок~\eqref{eq:latency} читається детекторами напрямку & Ілюзія працює й без кольору; пари сусідніх елементів породжують її поодинці; нейрони кори мавпи, вибіркові до напрямку, дають спрямовану відповідь на ці нерухомі пари. У людини обертання запускають мікрострибки очей і кліпання & \cite{Conway2005,OteroMillan2012} & виміряно \\
16 & Двозначні картинки: взаємне гальмування, втома й шум; зміни викликає головно шум & Модель конкурентних груп нейронів: зміни в ній викликає \emph{шум}, без шуму їх немає; вона пояснює зростання частоти змін із силою стимулу. Втома тут відіграє меншу роль & \cite{MorenoBote2007} & гіпотеза \\
17 & Частина ілюзій --- наслідок задачі, якої вчиться машина & Мережа, навчена передбачати наступний кадр відео, «бачить» обертання нерухомих кілець і не бачить його на контрольних картинках; мережі для очищення й підвищення різкості зображень піддаються ілюзіям кольору та яскравості, як людина & \cite{Watanabe2018,GomezVilla2020} & виміряно \\
\bottomrule
\end{longtable}}
